% Addendum to PAPER_B_MANUSCRIPT.tex: results the paper no longer carries, kept with
% their proofs.  Not replicated in second_order_sequence_stats_numbered_eqs.tex.
\documentclass[11pt,a4paper]{article}
\usepackage[margin=2.5cm]{geometry}
\usepackage{amsmath,amssymb,amsthm}

\newtheorem{proposition}{Proposition}

\newcommand{\PB}{P^{\mathrm{B}}}
\newcommand{\PJ}{P^{\mathrm{J}}}
\newcommand{\assoc}{\operatorname{assoc}}
\newcommand{\bigO}{\mathcal{O}}

\title{Addendum to ``Order-Sensitivity of Belief and Decision Statistics under Jeffrey Conditioning''}
\author{Anurag Srivastava}
\date{}

\begin{document}
\maketitle

This addendum keeps results that earlier drafts of the paper stated and the paper no
longer carries, each with its proof, so that they can be cited or restored. The setting
and the notation are the paper's, and each proposition keeps the name it had there.

\section{The adoption weight read from the ratings}

Proposition ADJ was removed from Section 6 of the paper on 3 October 2026. The paper keeps
its first identity as a sentence, since that identity is the partial-adoption rule read
backwards, and Proposition LAD(i) of the paper carries the gap at $c=0$. The second
identity, for the marginal read first, and the coefficient claim that follows from it are
stated only here.

Under partial adoption, after the first Jeffrey step the response to the second cue is a
Jeffrey step on its partition to the target $(1-\omega)\,m+\omega\,r_1$, where $m$ is the
marginal the first step left and $0\le\omega\le1$ is the adoption weight. Write
$P^{\omega}_{AB}$ and $P^{\omega}_{BA}$ for the two sequences under this rule; $\omega=1$
recovers $\PJ_{AB}$ and $\PJ_{BA}$, and $\omega=0$ ignores the second cue. Let $P^{A}$ and
$P^{B}$ denote the belief after the $A$-cue alone and after the $B$-cue alone.

\renewcommand{\theproposition}{ADJ}%
\begin{proposition}[adoption weight]\label{prop:ADJ}
For every $c$ and every $\omega$, in sequence $AB$,
\begin{align*}
P^{\omega}_{AB}(B{=}1)-r_1 &=(1-\omega)\bigl[P^{A}(B{=}1)-r_1\bigr],\\
P^{\omega}_{AB}(A{=}1)-q_1 &=\omega\,\bigl[\PJ_{AB}(A{=}1)-q_1\bigr],\\
P^{\omega}_{AB}(A{=}1)-P^{\omega}_{BA}(A{=}1)
  &=\omega\,\bigl[\PJ_{AB}(A{=}1)-q_1\bigr]+(1-\omega)\bigl[q_1-P^{B}(A{=}1)\bigr].
\end{align*}
Hence the last-read marginal equals its delivered credence identically in $c$
exactly when $\omega=1$, the first-read marginal does so exactly when
$\omega=0$, and for $0<\omega<1$ neither does. The first-order coefficient in
$c$ of the first-read marginal is $\omega$ times the benchmark's coefficient
for the same marginal. At $c=0$ the third line equals $(1-\omega)(\alpha-q_0)$,
so a partially adopting evaluator shows a sequence effect at independence, where
the fully adopting evaluator (Proposition IMM of the paper) and the benchmark
show none. Whenever $P^{A}(B{=}1)\neq r_1$, the first line gives the weight from
three marginals of one reading group, for every $c$,
\[
  \omega=\frac{P^{A}(B{=}1)-P^{\omega}_{AB}(B{=}1)}{P^{A}(B{=}1)-r_1}.
\]
\end{proposition}

\begin{proof}[Proof]

With $Q$ the belief after the first step and $t=(t_0,t_1)$,
$t_1=(1-\omega)\,Q(B{=}1)+\omega r_1$, the target of the second, column $j$
is multiplied by $t_j/Q(B{=}j)$, so that
\[
\begin{aligned}
  P(B{=}1)&=t_1,\\
  P(A{=}1)&=\sum_j Q(A{=}1,B{=}j)\,\frac{t_j}{Q(B{=}j)}\\
  &=(1-\omega)\,Q(A{=}1)+\omega\sum_j Q(A{=}1,B{=}j)\,\frac{r_j}{Q(B{=}j)}.
\end{aligned}
\]
In sequence $AB$, $Q=P^{A}$, $Q(A{=}1)=q_1$, and the last sum is
$\PJ_{AB}(A{=}1)$, giving
\[
  P^{\omega}_{AB}(B{=}1)-r_1=(1-\omega)\bigl[P^{A}(B{=}1)-r_1\bigr],\qquad
  P^{\omega}_{AB}(A{=}1)-q_1=\omega\bigl[\PJ_{AB}(A{=}1)-q_1\bigr].
\]
In sequence $BA$, by symmetry,
$P^{\omega}_{BA}(A{=}1)=(1-\omega)\,P^{B}(A{=}1)+\omega q_1$; subtracting
gives the third line. Neither bracket vanishes identically, since
\[
  P^{A}(B{=}1)-r_1\Big|_{c=0}=r_0-\beta,\qquad
  \PJ_{AB}(A{=}1)-q_1=-c\,\frac{q_0(1-q_0)(r_0-\beta)}{\alpha\beta(1-\alpha)(1-\beta)}+\bigO(c^2),
\]
the second by Proposition DRF of the paper, whose proof also gives
$\PJ_{AB}(A{=}1)-\PB(A{=}1)=\bigO(c^2)$, hence the coefficient claim. At
$c=0$ a step leaves the other marginal unchanged, so
\[
  \PJ_{AB}(A{=}1)=q_1,\qquad P^{B}(A{=}1)=1-\alpha,\qquad
  P^{\omega}_{AB}(A{=}1)-P^{\omega}_{BA}(A{=}1)=(1-\omega)(\alpha-q_0).
\]
For $P^{A}(B{=}1)\neq r_1$ the first line gives
\[
  1-\omega=\frac{P^{\omega}_{AB}(B{=}1)-r_1}{P^{A}(B{=}1)-r_1},\qquad
  \omega=\frac{P^{A}(B{=}1)-P^{\omega}_{AB}(B{=}1)}{P^{A}(B{=}1)-r_1}.
\]
\end{proof}

In Example~1 of the paper, an evaluator who reads the credential
first rates trustworthiness at $.56$ before the letter, the letter alone delivers
$.30$, and a final rating of $.43$ gives $\omega=(.56-.43)/(.56-.30)=\tfrac12$.

{\raggedright\paragraph{Records.} \texttt{lean/JeffreyOrder/Anchoring.lean} (\texttt{dampedB\_mB1},
\texttt{dampedB\_deviation}, \texttt{routeDamped\_mA1\_deviation},
\texttt{orderEffect\_damped\_at\_indep}); \texttt{sympy/verify\_ladder.py}.\par}

\end{document}
